\newpage

\section{Conclusiones}

El presente trabajo permitió establecer las bases conceptuales, técnicas y metodológicas para el desarrollo de un prototipo de sistema móvil orientado a la detección de indicios de infestación por insectos descortezadores en árboles, mediante el uso de visión por computadora y aprendizaje profundo. A partir del análisis del problema, se identificó que la detección de este tipo de plagas es una necesidad relevante para la conservación del arbolado urbano, debido a que los métodos tradicionales de inspección visual pueden resultar limitados ante la extensión de las áreas verdes, la falta de personal especializado y la rapidez con la que puede avanzar el deterioro de los árboles afectados.

La selección del Bosque de Chapultepec como caso de estudio permitió delimitar un escenario pertinente para el proyecto, al tratarse de un área verde urbana de alta importancia ecológica, social y cultural, en la que se ha documentado la presencia de afectaciones asociadas a insectos descortezadores. Asimismo, la caracterización de los hospederos, de los agentes causantes del daño y de las señales visuales presentes en el follaje y la corteza permitió definir los elementos observables que pueden ser aprovechados por un sistema computacional para apoyar la identificación de posibles casos de infestación.

Desde el punto de vista computacional, se concluye que la visión por computadora representa una alternativa viable para complementar las labores de monitoreo forestal, ya que permite analizar imágenes recolectadas en campo y extraer patrones visuales relacionados con el estado sanitario del arbolado. En este sentido, el enfoque de aprendizaje profundo resulta adecuado para el proyecto, debido a su capacidad para aprender características visuales complejas a partir de datos etiquetados, especialmente en tareas de clasificación de imágenes.

El análisis comparativo de arquitecturas permitió identificar modelos de referencia adecuados para el desarrollo del sistema, considerando criterios como precisión, eficiencia computacional y posibilidad de integración en dispositivos móviles. De igual forma, la definición del conjunto de datos, las etapas de preprocesamiento y la organización de imágenes para entrenamiento, validación y prueba constituyen elementos fundamentales para el desempeño del modelo de clasificación propuesto.

En cuanto al diseño del sistema, se estableció una arquitectura compuesta por una aplicación móvil para la captura y análisis de imágenes en campo, una plataforma web para la consulta y gestión de registros, y una infraestructura de soporte para el almacenamiento, procesamiento y seguimiento de la información generada. La definición de requerimientos funcionales y no funcionales, reglas de negocio, actores, casos de uso, flujos de operación y diseño de base de datos permitió estructurar una solución integral orientada tanto al análisis técnico como a la futura implementación del prototipo.

Se concluye que la propuesta es técnicamente viable como herramienta de apoyo para la detección y el seguimiento de árboles con posibles indicios de infestación. No obstante, su desempeño dependerá directamente de la calidad, cantidad y diversidad del conjunto de imágenes recolectadas, así como de la correcta validación de los resultados con especialistas en el área forestal. Factores como las condiciones de iluminación natural, el ángulo de captura, la variabilidad entre especies arbóreas y la similitud visual con otros tipos de daño representan limitaciones que deberán atenderse durante la fase de implementación.

Los resultados obtenidos en esta primera etapa sientan las bases necesarias para que en el Trabajo Terminal II se lleve a cabo la implementación, integración y validación del prototipo, abarcando el entrenamiento definitivo del modelo, el desarrollo de la aplicación móvil y la plataforma web, el despliegue de la infraestructura en la nube y la ejecución del plan de pruebas.